\section{How the Models Are Built and Kept Honest}
\label{sec:modelling}

Table~\ref{tab:scope} states what the artifact contains. The figures in it are
counted from the files themselves at build time, so the description cannot
drift away from the thing it describes.

\begin{table}[t]
\centering
\caption{Verification scope, measured from the repository rather than
described. ``Properties'' counts rows of the result matrix
(Table~\ref{tab:results}), not queries: a model file may carry several.}
\label{tab:scope}
\footnotesize
\begin{tabular}{@{}lrrr@{}}
\toprule
Component & Files & Lines & Properties \\
\midrule
Shared ProVerif core (\texttt{lib/}) & 3 & 349 & --- \\
ProVerif query models & 11 & 302 & 12 \\
Tamarin theory & 1 & 236 & 6 \\
Wire-format layout (\texttt{wire.py}) & 1 & 201 & 25 fields \\
\bottomrule
\end{tabular}

\end{table}

\subsection{One transcription, two provers}

The project's central methodological commitment is that any difference between
the two provers' outcomes should be attributable to the provers, not to the
modelling. The handshake is therefore written once, in
\texttt{lib/wireguard\_core.pvl}, and every ProVerif query file includes that
library rather than restating the protocol. The Tamarin model is transcribed
term for term against it, and Section~\ref{sec:discussion} records the two
places where the transcriptions could not be made identical, together with the
reason in each case.

\subsection{Primitives, and one abstraction that decided the project}

\para{Diffie-Hellman with two symbols.} Diffie-Hellman is modelled with two
symbols rather than an iterated exponentiation:
\[
  \mathit{pk}(x), \qquad
  \mathit{dh}(P,y), \qquad
  \mathit{dh}(\mathit{pk}(x),y)=\mathit{dh}(\mathit{pk}(y),x).
\]
Two consequences matter. Because \(\mathit{pk}\) is the only constructor of
the group type, every group element the attacker can supply is \(g^t\) for
some term \(t\) it can build---faithful for a prime-order group, where every
X25519 output is of that form. Because \(\mathit{dh}\) returns key material
rather than a group element, a Diffie-Hellman result cannot re-enter as a
group element, so term depth is bounded by construction.

\para{Why the obvious formulation fails.} This was not a stylistic choice. The
natural formulation \(\mathit{exp}(G,\mathit{exponent}):G\) lets the attacker
build towers \(\mathit{exp}(\mathit{exp}(\mathit{exp}(g,a),b),c)\) that the
single commuting equation does not normalise. With that formulation,
saturation on the secrecy query exhausted the memory budget after roughly
\(1.2\times10^4\) clauses and never terminated; with the two-symbol
formulation the same query is proved well inside the budget
(Table~\ref{tab:results}). The episode is recorded because it decides whether
a symbolic model is usable at all, and it is rarely written down.

\para{A flat transcript hash.} Noise updates \(h \leftarrow
\mathrm{HASH}(h \mathbin\| f)\) after each field. Since \(h\) is never
transmitted and serves only as AEAD associated data, its sole requirement is
that it bind the transcript prefix injectively, so we model \(h\) after step
\(i\) as a distinct collision-resistant function of the entire prefix. Both
formulations bind the same data; the flat one keeps term depth constant. AEAD,
the HKDF chain and the keyed MAC are standard perfect primitives, with one
function symbol per HKDF output position.

\subsection{The attacker, and when it gains each capability}

\begin{figure}[t]
\centering
\resizebox{\columnwidth}{!}{\begin{tikzpicture}[x=1cm,y=1cm,font=\footnotesize]
  \fill[okgreen!12] (0,0) rectangle (4.2,0.72);
  \fill[badred!12]  (4.2,0) rectangle (8.4,0.72);
  \node at (2.1,0.36) {\bfseries phase 0 --- sessions run};
  \node at (6.3,0.36) {\bfseries phase 1 --- after the fact};
  \draw[badred,dashed,thick] (4.2,-1.42) -- (4.2,0.95);

  \node[anchor=north west,text width=3.9cm,black!60,inner sep=2pt,
        font=\scriptsize] at (0.05,-0.06)
    {Dolev--Yao network control. Every ephemeral public value travels in
     clear, so all of it can be recorded.};
  \node[anchor=north west,text width=3.9cm,badred,inner sep=2pt,
        font=\scriptsize] at (4.3,-0.06)
    {Static scalars released and a discrete-logarithm oracle exposed: every
     recorded public value inverts.};

  \draw[black!25] (0,-1.6) -- (8.4,-1.6);
  \node[anchor=west,black!55,font=\scriptsize\itshape] at (0,-1.82)
    {which properties depend on this staging};
  \node[anchor=west] at (0,-2.20) {P3\quad forward secrecy, statics released};
  \node[anchor=east,okgreen] at (8.4,-2.20) {\bfseries proved};
  \node[anchor=west] at (0,-2.52) {P4\quad post-quantum FS, psk secret};
  \node[anchor=east,okgreen] at (8.4,-2.52) {\bfseries proved};
  \node[anchor=west] at (0,-2.84) {P4c\ \ the same model, psk published};
  \node[anchor=east,badred] at (8.4,-2.84) {\bfseries attack found};
  \node[anchor=west] at (0,-3.16) {S6a/S6b\ \ KCI, responder static leaked};
  \node[anchor=east,black!55] at (8.4,-3.16) {\bfseries the gap};
\end{tikzpicture}
}
\caption{Compromise is staged, never built in. Phase~0 is the classical
Dolev-Yao attacker; phase~1 releases the static scalars and a
discrete-logarithm oracle, strictly after every phase-0 session. The
key-compromise queries sit outside this staging, because there the compromise
is part of the scenario rather than a sequel to it.}
\label{fig:compromise}
\end{figure}

The attacker is Dolev-Yao. Compromise is never a built-in capability: each
class of secret is revealed by an explicit action, so every query states which
compromises it tolerates, and a proof cannot silently depend on an assumption
nobody wrote down.

\para{Staging a quantum transition.} Post-quantum capability is expressed with
ProVerif's phase mechanism (Figure~\ref{fig:compromise}). Phase~0 is the
classical attacker. In phase~1 the attacker gains a discrete-logarithm oracle,
private in phase~0 and exposed only by a phase-1 process. Since every
ephemeral public value travels in clear, an attacker that recorded phase-0
traffic can invert all of it once phase~1 begins. This is the standard
formalisation of harvest-now-decrypt-later, and it is what makes post-quantum
forward secrecy a property rather than a slogan.

\para{Compromise as part of the scenario.} Key-compromise impersonation needs
the opposite arrangement: the responder's static key is in the attacker's
hands \emph{while} the session runs, not afterwards. That is an ordering
constraint over individual events rather than a global staging, which is why
those queries are posed to Tamarin with action facts and why the two provers
do not divide this work symmetrically.

\subsection{From wire format to model terms}

A formal model earns its authority from the fidelity of its correspondence to
the artifact it claims to describe. We therefore maintain a field-by-field
mapping from the wire format to terms in the models, and make it
machine-checkable rather than a promise.

\para{Offsets are derived, never written.} \texttt{tools/wire.py} is the
single source of truth for the layout. Offsets are accumulated from the field
lengths, so a corrected length cannot leave a stale offset behind it, and the
computed totals are checked against the specification figures of
\wgInitBytes, \wgRespBytes, 64 and 16\,B. Tables~\ref{tab:wireinit}
and~\ref{tab:wireresp} are rendered from that same structure, which is also
what \texttt{scripts/verify\_wire\_table.py} cross-checks the models against.
Of \wgFields{} fields across the four message types, \wgFieldsModelled{} carry
a model term and \wgFieldsUnmodelled{} deliberately do not.

\begin{table}[t]
\centering
\caption{Handshake initiation, \wgInitBytes\,B. Offsets are derived from
lengths; ``model term'' names the corresponding term in the shared ProVerif
core, and a dash marks a field deliberately outside the models.}
\label{tab:wireinit}
\footnotesize
\begin{tabular}{@{}rrll@{}}
\toprule
Off & Len & Field & Model term \\
\midrule
0 & 1 & \texttt{message\_type} & \textcolor{black!45}{---} \\
1 & 3 & \texttt{reserved} & \textcolor{black!45}{---} \\
4 & 4 & \texttt{sender\_index} & \textcolor{black!45}{---} \\
8 & 32 & \texttt{unencrypted\_ephemeral} & $E_i=\mathit{pk}(e_i)$ \\
40 & 48 & \texttt{encrypted\_static} & $\mathrm{aead}(k_1,0,S_i,h_1)$ \\
88 & 28 & \texttt{encrypted\_timestamp} & $\mathrm{aead}(k_2,0,\mathit{ts},h_2)$ \\
116 & 16 & \texttt{mac1} & $\mathrm{mac}(H(L\|S_r),\cdot)$ \\
132 & 16 & \texttt{mac2} & \textcolor{black!45}{---} \\
\bottomrule
\end{tabular}

\end{table}

\begin{table}[t]
\centering
\caption{Handshake response, \wgRespBytes\,B. Note that \texttt{mac1} here is
keyed by the \emph{initiator's} static public value, which is what makes
identity hiding conditional.}
\label{tab:wireresp}
\footnotesize
\begin{tabular}{@{}rrll@{}}
\toprule
Off & Len & Field & Model term \\
\midrule
0 & 1 & \texttt{message\_type} & \textcolor{black!45}{---} \\
1 & 3 & \texttt{reserved} & \textcolor{black!45}{---} \\
4 & 4 & \texttt{sender\_index} & \textcolor{black!45}{---} \\
8 & 4 & \texttt{receiver\_index} & \textcolor{black!45}{---} \\
12 & 32 & \texttt{unencrypted\_ephemeral} & $E_r=\mathit{pk}(e_r)$ \\
44 & 16 & \texttt{encrypted\_nothing} & $\mathrm{aead}(k_3,0,\varepsilon,h_5)$ \\
60 & 16 & \texttt{mac1} & $\mathrm{mac}(H(L\|S_i),\cdot)$ \\
76 & 16 & \texttt{mac2} & \textcolor{black!45}{---} \\
\bottomrule
\end{tabular}

\end{table}

\para{What is not modelled, and why.} Four field classes are deliberately
absent, and stating the reason for each is part of the correspondence rather
than an apology for it. \emph{Message type and reserved bytes} serve
demultiplexing; the models use distinct channels instead. \emph{Session
indices} are chosen freely by the sender and carry no security burden, but
they are also what makes endpoint roaming possible---a WireGuard engineering
advantage with no counterpart in an abstract Noise model. \emph{\texttt{mac2}
and the cookie mechanism} are load-shedding devices, and the Dolev-Yao
attacker has no notion of computational cost, so a symbolic model cannot
express the property they provide. \emph{The data-plane counter and its
anti-replay window} lie outside the handshake entirely.

\para{An argued property, not a verified one.} The cookie mechanism deserves
the distinction made explicitly. \texttt{mac1} is keyed by
\(\mathrm{HASH}(\mathtt{LABEL\_MAC1} \mathbin\| S_r)\), so an attacker
ignorant of the responder's static public value cannot produce a message that
survives the first check, and cannot force a Diffie-Hellman operation. That is
a genuine property, and it is genuinely outside the model. We argue it here
rather than claim the models establish it.

\para{Conditional identity hiding.} One field deserves separate mention. As
Table~\ref{tab:wireresp} records, the \texttt{mac1} of the \emph{response} is
keyed by the \emph{initiator's} static public value. An observer holding a
candidate public key can therefore test whether a given response is addressed
to its holder. WireGuard's identity hiding is accordingly conditional: it
protects the initiator's identity from an observer who does not already hold
the candidate keys, and not from one who does. The distinction is visible only
because \texttt{mac1} was modelled explicitly rather than abstracted away.

\subsection{Budgets, expectations, and generated numbers}

\para{Every run is bounded.} Every invocation carries an explicit
address-space cap and wall-clock timeout---\wgMemCap\,GB and \wgBudget\,s on a
single core, a budget chosen to match what remains available on an 8\,GB
laptop with an editor open. The cap matters more than its value: without one,
a runaway query consumes the machine and produces a run that is slow,
irreproducible and hostile to anything else running. With one, a query either
finishes within budget or fails deterministically, which is a reportable
outcome.

\para{Verification as a regression suite.}
\texttt{expectations.yaml} declares the expected verdict for each of
\wgExpectations{} queries, and \texttt{make verify} fails on any mismatch.
Outcomes that are attacks are expectations too: the control model of
Section~\ref{sec:results} is \emph{supposed} to be refuted, and a run in which
it succeeded would indicate that the model had drifted, not that the protocol
had improved.

\para{No number in this paper is typed.} Result tables, figures and every
quantity this text quotes are rendered from \texttt{results/results.json},
\texttt{tools/wire.py} and \texttt{tools/pqcost.py} when the document is
built. An earlier revision of this report was internally consistent and
externally wrong---the tables were generated and correct while the prose
quoted timings from a run nobody could identify. Generating the tables was not
enough, because prose is not data; the numbers in the sentences are now macros
too.
